% TOP_PROBLEM: {"id":30007637,"problem_number":"LOCAL-30007637","title":"Wages versus job amenities","category_id":21,"category":{"id":21,"name":"economics","display_name":"Economics","description":"Open economic theory statements, published research agendas, and proposed scoped specifications; question provenance and historical priority are qualified per record.","slug":"economics","order_index":21,"created_at":"2026-10-08T00:00:00Z"},"status":"research_question","external_url":null,"legacy_ids":[],"published":false,"statement":"In the explicitly specified setting: How should earnings inequality be adjusted for heterogeneous preferences over flexibility, security, hours, and working conditions when labor markets are imperfectly competitive? The mathematical statement above fixes the target, assumptions and scope of this catalogue formulation.","economics":{"original_id":126,"field":"Labor markets and work","kind":"identification","formalization_basis":"adapted_research_question","source_status":"research_agenda","priority_status":"earliest_found","priority_note":"Earliest directly inspected qualifying passage in this audit; earlier origins were not established. A review or research-agenda statement does not imply original priority.","earliest_verified_reference":{"authors":["Alexandre Mas"],"title":"Non-Wage Amenities","year":2025,"url":"https://www.rfberlin.com/wp-content/uploads/2025/08/25051.pdf","doi":"10.3386/w33643","locator":"Section 8.2, printed p. 63; conclusion, printed pp. 75–76","evidence_summary":"Estimating preferences from job choices requires feasible consideration sets and an offer process; omitted productivity biases compensating-differential estimates."},"current_reference":{"authors":["Alexandre Mas"],"title":"Non-Wage Amenities","year":2025,"url":"https://www.rfberlin.com/wp-content/uploads/2025/08/25051.pdf","doi":"10.3386/w33643","locator":"Section 8.2, printed p. 63; conclusion, printed pp. 75–76","evidence_summary":"Estimating preferences from job choices requires feasible consideration sets and an offer process; omitted productivity biases compensating-differential estimates."},"formalization_note":"The cited passage establishes a research direction, not this exact mathematical statement. The notation, population, policy menu, assumptions, and estimands are an original adaptation for this catalogue; no universal unsolved theorem or source priority is claimed."},"statusQualification":"Published question or research agenda verified at the cited locator. The mathematical specification is adapted for this catalogue and includes additional assumptions; source publication does not establish first-ever priority or certify this exact model remains unsolved. The cited passage establishes a research direction, not this exact mathematical statement. The notation, population, policy menu, assumptions, and estimands are an original adaptation for this catalogue; no universal unsolved theorem or source priority is claimed.","statusReviewedAt":"2026-10-08"}
% FIRST_POSTED: 2026-10-08
% ENTRY_KIND: statement_only
% !TeX program = lualatex
\documentclass[11pt]{article}
\usepackage{fontspec}
\usepackage[a4paper,margin=25mm]{geometry}
\usepackage{amsmath,amssymb,mathtools}
\usepackage{xurl}
\usepackage[hidelinks]{hyperref}
\newcommand{\sref}[2]{\hyperlink{source:#1:#2}{[S#2]}}
\setlength{\emergencystretch}{3em}
\begin{document}
% problemId:30007637
\section*{Wages versus job amenities}\label{30007637}
\subsection{Definitions and mathematical statement}
Fix a worker population and an observable set of jobs. Job \(j\) supplies wage \(w_j\) and amenity vector \(a_j\), including scheduling, security, location, and workplace conditions. Worker \(i\) has utility \(u_i(w,a)\), strictly increasing in wage, and a feasible consideration set \(C_i\). For a declared common amenity bundle \(a_0\), define wage-equivalent job value \(e_{ij}\) by \(u_i(e_{ij},a_0)=u_i(w_j,a_j)\), when this inverse exists. Let \(j(i)\) denote actual employment. The target is the distribution of \(e_{i,j(i)}\), together with its inequality relative to the distribution of observed wages. Identify or bound the required preferences using credible offer variation, stated choices validated against behavior, and information about \(C_i\). Separate compensation for amenities from firm productivity, rent sharing, and endogenous matching. Report how conclusions depend on the reference bundle and heterogeneous preferences; interpersonal utility comparisons need additional declared normalization. A policy counterfactual equalizing amenities must additionally model wage and job reallocation. The unresolved measurement task is to quantify these adjustments reliably in a specified population, not to infer their sign from competitive hedonic regressions alone.

\par\textbf{Formulation and scope.} The cited passage establishes a research direction, not this exact mathematical statement. The notation, population, policy menu, assumptions, and estimands are an original adaptation for this catalogue; no universal unsolved theorem or source priority is claimed.

\par\textbf{Open-question provenance.} An explicit question or research agenda was verified at the source locator. This does not by itself certify that every later version remains unresolved \sref{30007637}{1}.

\par\textbf{Historical priority.} Earliest directly inspected qualifying passage in this audit; earlier origins were not established. A review or research-agenda statement does not imply original priority.

\subsection{Short English statement}
In the explicitly specified setting: How should earnings inequality be adjusted for heterogeneous preferences over flexibility, security, hours, and working conditions when labor markets are imperfectly competitive? The mathematical statement above fixes the target, assumptions and scope of this catalogue formulation.

\subsection{Sources}
\begin{itemize}\raggedright
\item[\textnormal{[S1]}]\hypertarget{source:30007637:1}{} Alexandre Mas. 2025. Non-Wage Amenities. Section 8.2, printed p. 63; conclusion, printed pp. 75–76.
\url{https://www.rfberlin.com/wp-content/uploads/2025/08/25051.pdf}.
DOI: 10.3386/w33643.
{\small\textit{Earliest verified question/agenda reference; historical priority qualified below. Estimating preferences from job choices requires feasible consideration sets and an offer process; omitted productivity biases compensating-differential estimates.}}
\end{itemize}
\end{document}
